\section{到达时刻、指数等待与条件均匀性}
\label{sec:06-Random-Processes/02-Counting-Processes/03}

Poisson 过程有两种等价描述。上一节从计数增量出发——``一段时间内发生多少次''。本节从间隔出发——``下一个事件还要等多久''。两种描述回答不同的问题，却能互相翻译。

令第 \(r\) 个到达时刻为

\[
T_r=\inf\{t:N(t)\ge r\},
\]

并令相邻间隔

\[
W_1=T_1,\qquad
W_r=T_r-T_{r-1}.
\]

\subsection{第一次等待：为什么是指数}

你想知道第一个事件什么时候来。事件 \(\{T_1>t\}\)——直到时刻 \(t\) 还没有事件——等价于计数为零：

\[
\{T_1> t\}=\{N(t)=0\}.
\]

而 Poisson 过程在长度为 \(t\) 的区间内事件数为零的概率是

\[
P(N(t)=0)=e^{-\lambda t}.
\]

所以

\[
P(T_1> t)=e^{-\lambda t},
\]

这正是 Exponential 分布的生存函数。于是

\[
T_1\sim\mathrm{Exponential}(\lambda).
\]

这与离散时间的 Geometric 完全平行：\(P(T_1=k)=(1-p)^{k-1}p\) 是``连续失败后一次成功''，\(e^{-\lambda t}\lambda\,dt\) 是``持续无事件后一次事件''。二项极限让 Geometric 变成了 Exponential——将格点间距缩为零，离散等待变成了连续等待。

\subsection{间隔独立同分布：过程在每个到达时刻重启}

一次事件发生后，剩下的过程和从头开始没有区别——因为你把那一次事件之前的计数减掉，独立增量性质告诉你剩余部分仍然是速率为 \(\lambda\) 的 Poisson 过程，与过去独立。

这个直觉需要小心论证：到达时刻 \(T_r\) 是随机的，而独立增量性质严格说只对固定时间点成立。需要更强的强 Markov 性质——在随机停止时间处过程仍然``重启''。成立，但细节属于测度论范畴。

结论是：

\[
W_1,W_2,\ldots
\overset{\text{iid}}{\sim}
\mathrm{Exponential}(\lambda).
\]

Exponential 的无记忆性

\[
P(W> s+t\mid W> s)=P(W> t)
\]

恰好表达了``已经等了多久都不会改变剩余等待分布''。这既是 Markov 性质在间隔层面的体现，也说明指数分布是连续分布中唯一具有无记忆性的分布——如果等待时间有记忆，就一定不是指数分布。

\subsection{第 r 个到达时刻：独立指数和}

既然 \(T_r\) 是 \(r\) 个独立 Exponential 变量之和：

\[
T_r=W_1+\cdots+W_r,
\]

而独立 Exponential 之和服从 Gamma 分布（卷积或比较矩母函数均可验证），所以

\[
T_r\sim\mathrm{Gamma}(r,\text{rate }\lambda).
\]

密度为

\[
f_{T_r}(t)
=\frac{\lambda^r t^{r-1}e^{-\lambda t}}
{(r-1)!},
\qquad t>0.
\]

Gamma\((r,\lambda)\) 在排队论中也称为 Erlang 分布——Erlang 在电话交换研究中引入了这个模型，用来描述 \(r\) 个接线员依次处理呼叫的总耗时。

同一个结论也可以从计数视角得到：事件 \(\{T_r\le t\}=\{N(t)\ge r\}\) 给出了 \(T_r\) 的累积分布与 Poisson 尾概率的关系。这和 Bernoulli 过程中 \(P(T_r\le n)=P(N_n\ge r)\) 是同一种对偶——一种遍布计数过程的通用语言，不限于 Poisson。

\subsection{条件均匀性：给定总数后，到达位置反而均匀}

现在出现一个反直觉的结论。

条件于截止时刻 \(t\) 的事件总数

\[
N(t)=n,
\]

区间 \((0,t)\) 中的 \(n\) 个无序到达位置，与 \(n\) 个独立 Uniform\((0,t)\) 样本同分布。按时间排序后，这些位置就是均匀次序统计量。

具体说：如果一小时内恰有一次到达，它发生时刻在这一小时内的任何位置等可能。两小时内有三次到达，这三次的到达时间就像在 \((0,2)\) 区间内随机撒了三个点然后排序。

这似乎和无条件间隔的 Exponential 分布矛盾——指数分布的概率密度在零附近最大，不应该``更可能早到''吗？

不矛盾。两个陈述条件不同：
\begin{itemize}
\tightlist
\item
  无条件下，第一次到达倾向于更早——因为如果晚了，中间可能已经发生了多次事件，不再``只有一次''；
\item
  条件于``恰好一次''，你排除了``发生多次''的可能。在唯一一次到达的条件下，它落在任何位置的密度均匀。
\end{itemize}

给定 \(N(t)=1\)，唯一到达落在前半段的概率为 \(1/2\)。但不条件于总数，第一次到达更可能靠近零——早到了后面还可以再来，但来晚了就没有``第一次''的位置了。

\subsection{次序统计量的分布}

进一步，给定 \(N(t)=n\)，第 \(r\) 个到达的时间比例

\[
\frac{T_r}{t}
\]

服从

\[
\mathrm{Beta}(r,n-r+1),
\qquad r=1,\ldots,n.
\]

Beta 分布是次序统计量的标准分布——\(n\) 个独立均匀样本的第 \(r\) 次序统计量刚好是 Beta\((r,n-r+1)\)。这也印证了条件均匀性。

条件均值为

\[
\E[T_r\mid N(t)=n]
=\frac{r}{n+1}t.
\]

分母是 \(n+1\) 而不是 \(n\)：到达点平均把区间等分成 \(n+1\) 段。有三个人在一小时内到达，他们的到达时刻大约分别在 15 分钟、30 分钟和 45 分钟——均匀铺开，而不是聚集在某个时间段。

\subsection{无记忆性：方便的工具，苛刻的假设}

指数分布的无记忆性让分析极度简化。如果服务时间服从指数分布，你只需要记录系统当前有多少顾客，不需要记录每个人已经服务了多久——因为剩余服务时间的分布和刚进入系统时完全一样。这正是连续时间 Markov 链和 M/M/1 队列可以用低维状态空间描述的根本原因。

但现实中：
\begin{itemize}
\tightlist
\item
  机器的故障率可能随年龄上升（老化）——剩余寿命依赖已使用时间；
\item
  服务可能包含多个固定阶段——完成第一阶段后剩余时间不再是指数分布；
\item
  人的耐心有限——等待越久，放弃的概率越大。
\end{itemize}

此时需要扩充状态空间（记录年龄、阶段）或转向更新过程、半 Markov 模型。放弃指数间隔后如何系统分析，正是后续 Markov 链单元要处理的问题——那里你会看到如何用状态分解来替代无记忆假设。
